% GENERATED FILE. Edit canonical lessons, then run npm run generate:books.
% canonical-source-hash: fnv1a64:38e0d4ddd5b4a92a
% canonical-lessons: MR-C487-catka-basne, MR-C487-galu, MR-C487-pural, MR-C487-bhajne, MR-C487-khajavne

\chapter{Burns and Rashes}
\label{ch:mr-burns-and-rashes}
\hlchaptermodality{487}

\begin{chapteropening}
\textbf{By the end of this chapter:} \emph{I can talk about a burn, a boil, a rash and an itch.}

Complete the last lesson of chapter 487: I can talk about a burn, a boil, a rash and an itch.
\end{chapteropening}

\section[\texorpdfstring{\mr{चटका} \mr{बसणे} (caṭkā basṇe)}{चटका बसणे (caṭkā basṇe)}]{\mr{चटका} \mr{बसणे} (caṭkā basṇe) — to get a burn}

\label{lesson:MR-C487-catka-basne}



\begin{quote}
Before the new one: say the Marathi for to sprain, then the Marathi for to swell.
\end{quote}

\subsection*{You'll want to know: \mr{चटका} \mr{बसणे}}
\textbf{\mr{चटका} \mr{बसणे}} — \emph{caṭkā basṇe} — \textquotedblleft{}to get a burn\textquotedblright{}.

\begin{quote}
\mr{गरम} \mr{तव्याला} \mr{हात} \mr{लागला} \mr{आणि} \mr{चटका} \mr{बसला}.
\end{quote}

\begin{quote}
\emph{garam tavyālā hāt lāglā āṇi caṭkā baslā.}
\end{quote}

\begin{quote}
My hand touched the hot griddle and I got burnt.
\end{quote}

Past (to him …): \mr{त्याला} \mr{चटका} \mr{बसला}, \emph{tyālā caṭkā baslā}.

Dative: \mr{मला} \mr{चटका} \mr{बसला}. Also a sudden pang of grief.

\begin{cousinweb}
\textbf{\mr{चटका} \mr{बसणे}} (\emph{caṭkā basṇe}) is a Marathi phrase built on \mr{चटका}, which CDIAL lists as Marathi \emph{caṭkā}, \textquotedblleft{}glow\textquotedblright{}, under \emph{*caṭa} (CDIAL 4570), with \mr{बसणे}.
\end{cousinweb}

\begin{grammarlens}[title={What you've built}]
One more word for describing what you lived and felt.
\end{grammarlens}

\subsection*{Your turn}
\begin{itemize}
\raggedright
  \item \emph{Say it:} \emph{caṭkā basṇe}
  \item \emph{Say it:} \emph{caṭkā basṇe}, once more
  \item \emph{Say it:} \emph{garam tavyālā hāt lāglā āṇi caṭkā baslā.} — My hand touched the hot griddle and I got burnt.
  \item \emph{Recall:} say the Marathi for to sprain, then the Marathi for to swell, then say \emph{caṭkā basṇe} again
\end{itemize}

\begin{tcolorbox}[breakable,colback=teal!4,colframe=teal!35!black,title={Before you move on}]
What does \mr{चटका} \mr{बसणे} mean? (\textquotedblleft{}To get a burn\textquotedblright{}.) Say it once more.
\end{tcolorbox}
\section[\texorpdfstring{\mr{गळू} (gaḷū)}{गळू (gaḷū)}]{\mr{गळू} (gaḷū) — an abscess; a boil}

\label{lesson:MR-C487-galu}



\begin{quote}
Before the new one: say the Marathi for to swell, then the Marathi for to get a burn.
\end{quote}

\subsection*{You'll want to know: \mr{गळू}}
\textbf{\mr{गळू}} — \emph{gaḷū} — \textquotedblleft{}an abscess; a boil\textquotedblright{}.

\begin{quote}
\mr{पाठीवरचं} \mr{गळू} \mr{आता} \mr{फुटलं}.
\end{quote}

\begin{quote}
\emph{pāṭhīvarcaṁ gaḷū ātā phuṭlaṁ.}
\end{quote}

\begin{quote}
The boil on the back has burst now.
\end{quote}

Neuter.

Neuter; \mr{गळू} \mr{होणे}, to get a boil.

\begin{cousinweb}
\textbf{\mr{गळू}} (\emph{gaḷū}) is a tadbhava: CDIAL lists Marathi \emph{gaḷū̃}, \textquotedblleft{}boil\textquotedblright{}, under Sanskrit \emph{gaḍu} (CDIAL 3977).
\end{cousinweb}

\begin{grammarlens}[title={What you've built}]
One more word for describing what you lived and felt.
\end{grammarlens}

\subsection*{Your turn}
\begin{itemize}
\raggedright
  \item \emph{Say it:} \emph{gaḷū}
  \item \emph{Say it:} \emph{gaḷū}, once more
  \item \emph{Say it:} \emph{pāṭhīvarcaṁ gaḷū ātā phuṭlaṁ.} — The boil on the back has burst now.
  \item \emph{Recall:} say the Marathi for to swell, then the Marathi for to get a burn, then say \emph{gaḷū} again
\end{itemize}

\begin{tcolorbox}[breakable,colback=teal!4,colframe=teal!35!black,title={Before you move on}]
What does \mr{गळू} mean? (\textquotedblleft{}An abscess; a boil\textquotedblright{}.) Say it once more.
\end{tcolorbox}
\section[\texorpdfstring{\mr{पुरळ} (puraḷ)}{पुरळ (puraḷ)}]{\mr{पुरळ} (puraḷ) — a rash; spots}

\label{lesson:MR-C487-pural}



\begin{quote}
Before the new one: say the Marathi for to get a burn, then the Marathi for an abscess.
\end{quote}

\subsection*{You'll want to know: \mr{पुरळ}}
\textbf{\mr{पुरळ}} — \emph{puraḷ} — \textquotedblleft{}a rash; spots\textquotedblright{}.

\begin{quote}
\mr{उन्हामुळे} \mr{मुलाच्या} \mr{पाठीवर} \mr{पुरळ} \mr{आलं}.
\end{quote}

\begin{quote}
\emph{unhāmuḷe mulācyā pāṭhīvar puraḷ ālaṁ.}
\end{quote}

\begin{quote}
The sun brought out a rash on the boy's back.
\end{quote}

Neuter.

Neuter; \mr{पुरळ} \mr{येणे}, to break out in a rash.

\begin{cousinweb}
\textbf{\mr{पुरळ}} (\emph{puraḷ}) is a Marathi noun whose origin is not settled.
\end{cousinweb}

\begin{grammarlens}[title={What you've built}]
One more word for describing what you lived and felt.
\end{grammarlens}

\subsection*{Your turn}
\begin{itemize}
\raggedright
  \item \emph{Say it:} \emph{puraḷ}
  \item \emph{Say it:} \emph{puraḷ}, once more
  \item \emph{Say it:} \emph{unhāmuḷe mulācyā pāṭhīvar puraḷ ālaṁ.} — The sun brought out a rash on the boy's back.
  \item \emph{Recall:} say the Marathi for to get a burn, then the Marathi for an abscess, then say \emph{puraḷ} again
\end{itemize}

\begin{tcolorbox}[breakable,colback=teal!4,colframe=teal!35!black,title={Before you move on}]
What does \mr{पुरळ} mean? (\textquotedblleft{}A rash; spots\textquotedblright{}.) Say it once more.
\end{tcolorbox}
\section[\texorpdfstring{\mr{भाजणे} (bhājṇe)}{भाजणे (bhājṇe)}]{\mr{भाजणे} (bhājṇe) — to roast; to get burnt}

\label{lesson:MR-C487-bhajne}



\begin{quote}
Before the new one: say the Marathi for an abscess, then the Marathi for a rash.
\end{quote}

\subsection*{You'll want to know: \mr{भाजणे}}
\textbf{\mr{भाजणे}} — \emph{bhājṇe} — \textquotedblleft{}to roast; to get burnt\textquotedblright{}.

\begin{quote}
\mr{स्वयंपाक} \mr{करताना} \mr{ती} \mr{भाजली}.
\end{quote}

\begin{quote}
\emph{svayampāk kartānā tī bhājlī.}
\end{quote}

\begin{quote}
She got burnt while cooking.
\end{quote}

Past (he …): \mr{तो} \mr{भाजला}, \emph{to bhājlā}.

Of people, to get burnt; of food, to roast: \mr{शेंगदाणे} \mr{भाजणे}.

\begin{cousinweb}
\textbf{\mr{भाजणे}} (\emph{bhājṇe}) is a tadbhava: CDIAL lists Marathi \emph{bhājṇẽ}, \textquotedblleft{}fries, parches\textquotedblright{}, under Sanskrit \emph{bhṛjjati} (CDIAL 9583).
\end{cousinweb}

\begin{grammarlens}[title={What you've built}]
One more word for describing what you lived and felt.
\end{grammarlens}

\subsection*{Your turn}
\begin{itemize}
\raggedright
  \item \emph{Say it:} \emph{bhājṇe}
  \item \emph{Say it:} \emph{bhājṇe}, once more
  \item \emph{Say it:} \emph{svayampāk kartānā tī bhājlī.} — She got burnt while cooking.
  \item \emph{Recall:} say the Marathi for an abscess, then the Marathi for a rash, then say \emph{bhājṇe} again
\end{itemize}

\begin{tcolorbox}[breakable,colback=teal!4,colframe=teal!35!black,title={Before you move on}]
What does \mr{भाजणे} mean? (\textquotedblleft{}To roast; to get burnt\textquotedblright{}.) Say it once more.
\end{tcolorbox}
\section[\texorpdfstring{\mr{खाजवणे} (khājavṇe)}{खाजवणे (khājavṇe)}]{\mr{खाजवणे} (khājavṇe) — to scratch (an itch)}

\label{lesson:MR-C487-khajavne}



\begin{quote}
Before the new one: say the Marathi for a rash, then the Marathi for to roast.
\end{quote}

\subsection*{You'll want to know: \mr{खाजवणे}}
\textbf{\mr{खाजवणे}} — \emph{khājavṇe} — \textquotedblleft{}to scratch (an itch)\textquotedblright{}.

\begin{quote}
\mr{डास} \mr{चावल्यावर} \mr{तो} \mr{पाय} \mr{खाजवत} \mr{होता}.
\end{quote}

\begin{quote}
\emph{ḍās cāvlyāvar to pāy khājavat hotā.}
\end{quote}

\begin{quote}
After the mosquito bit him he was scratching his leg.
\end{quote}

Past, with -\mr{ने} (he …): \mr{त्याने} \mr{पाय} \mr{खाजवला}, \emph{tyāne pāy khājavlā}.

\mr{खाज}, an itch: \mr{खाज} \mr{येणे}, to itch, with the dative.

\begin{cousinweb}
\textbf{\mr{खाजवणे}} (\emph{khājavṇe}) is a Marathi verb built on \mr{खाज}, \textquotedblleft{}itch\textquotedblright{}, which CDIAL lists under Sanskrit \emph{kharju}, \textquotedblleft{}itching, scab\textquotedblright{} (CDIAL 3827).
\end{cousinweb}

\begin{grammarlens}[title={What you've built}]
One more word for describing what you lived and felt.
\end{grammarlens}

\subsection*{Your turn}
\begin{itemize}
\raggedright
  \item \emph{Say it:} \emph{khājavṇe}
  \item \emph{Say it:} \emph{khājavṇe}, once more
  \item \emph{Say it:} \emph{ḍās cāvlyāvar to pāy khājavat hotā.} — After the mosquito bit him he was scratching his leg.
  \item \emph{Recall:} say the Marathi for a rash, then the Marathi for to roast, then say \emph{khājavṇe} again
\end{itemize}

\begin{tcolorbox}[breakable,colback=teal!4,colframe=teal!35!black,title={Before you move on}]
What does \mr{खाजवणे} mean? (\textquotedblleft{}To scratch (an itch)\textquotedblright{}.) Say it once more.
\end{tcolorbox}
